\documentclass[10pt,a4paper]{amsart}
\usepackage{amsmath,amssymb,amsthm,graphicx}
\usepackage{verbatim}
\newenvironment{reference}{\comment}{\endcomment}
%\newenvironment{reference}{}{}
\newenvironment{slogan}{\comment}{\endcomment}
\newenvironment{history}{\comment}{\endcomment}

% For commutative diagrams we use Xy-pic
\usepackage[all]{xy}

% We use 2cell for 2-commutative diagrams.
\xyoption{2cell}
% Theorem environments.
%
\theoremstyle{plain}
\newtheorem{theorem}[subsection]{Theorem}
\newtheorem{proposition}[subsection]{Proposition}
\newtheorem{lemma}[subsection]{Lemma}

\theoremstyle{definition}
\newtheorem{definition}[subsection]{Definition}
\newtheorem{example}[subsection]{Example}
\newtheorem{exercise}[subsection]{Exercise}
\newtheorem{situation}[subsection]{Situation}

\theoremstyle{remark}
\newtheorem{remark}[subsection]{Remark}
\newtheorem{remarks}[subsection]{Remarks}

\numberwithin{equation}{subsection}

% Macros
%
\def\lim{\mathop{\mathrm{lim}}\nolimits}
\def\colim{\mathop{\mathrm{colim}}\nolimits}
\def\Spec{\mathop{\mathrm{Spec}}}
\def\Hom{\mathop{\mathrm{Hom}}\nolimits}
\def\Ext{\mathop{\mathrm{Ext}}\nolimits}
\def\SheafHom{\mathop{\mathcal{H}\!\mathit{om}}\nolimits}
\def\SheafExt{\mathop{\mathcal{E}\!\mathit{xt}}\nolimits}
\def\Sch{\mathit{Sch}}
\def\Mor{\mathop{\mathrm{Mor}}\nolimits}
\def\Ob{\mathop{\mathrm{Ob}}\nolimits}
\def\Sh{\mathop{\mathit{Sh}}\nolimits}
\def\NL{\mathop{N\!L}\nolimits}
\def\CH{\mathop{\mathrm{CH}}\nolimits}
\def\proetale{{pro\text{-}\acute{e}tale}}
\def\etale{{\acute{e}tale}}
\def\QCoh{\mathit{QCoh}}
\def\Ker{\mathop{\mathrm{Ker}}}
\def\Im{\mathop{\mathrm{Im}}}
\def\Coker{\mathop{\mathrm{Coker}}}
\def\Coim{\mathop{\mathrm{Coim}}}

% Boxtimes
%
\DeclareMathSymbol{\boxtimes}{\mathbin}{AMSa}{"02}

%
% Macros for moduli stacks/spaces
%
\def\QCohstack{\mathcal{QC}\!\mathit{oh}}
\def\Cohstack{\mathcal{C}\!\mathit{oh}}
\def\Spacesstack{\mathcal{S}\!\mathit{paces}}
\def\Quotfunctor{\mathrm{Quot}}
\def\Hilbfunctor{\mathrm{Hilb}}
\def\Curvesstack{\mathcal{C}\!\mathit{urves}}
\def\Polarizedstack{\mathcal{P}\!\mathit{olarized}}
\def\Complexesstack{\mathcal{C}\!\mathit{omplexes}}
% \Pic is the operator that assigns to X its picard group, usage \Pic(X)
% \Picardstack_{X/B} denotes the Picard stack of X over B
% \Picardfunctor_{X/B} denotes the Picard functor of X over B
\def\Pic{\mathop{\mathrm{Pic}}\nolimits}
\def\Picardstack{\mathcal{P}\!\mathit{ic}}
\def\Picardfunctor{\mathrm{Pic}}
\def\Deformationcategory{\mathcal{D}\!\mathit{ef}}
\usepackage{cleveref}
\crefname{originalequation}{original equation}{original equations}
\Crefname{originalequation}{Original equation}{Original equations}
\makeatletter
\newlabel{equation-base-change}{{0A9J}{source}{}{sourcepointer-1}{}}
\newlabel{equation-base-change@cref}{{[originalequation][0][]0A9J}{[1][0][]source}{}{}{}}
\makeatother
\crefname{originalreference}{original reference}{original references}
\Crefname{originalreference}{Original reference}{Original references}
\makeatletter
\newlabel{cohomology-remark-base-change}{{08HY}{source}{}{sourcepointer-2}{}}
\newlabel{cohomology-remark-base-change@cref}{{[originalreference][0][]08HY}{[1][0][]source}{}{}{}}
\makeatother
\crefname{originalequation}{original equation}{original equations}
\Crefname{originalequation}{Original equation}{Original equations}
\makeatletter
\newlabel{equation-base-change-map}{{0AA6}{source}{}{sourcepointer-3}{}}
\newlabel{equation-base-change-map@cref}{{[originalequation][0][]0AA6}{[1][0][]source}{}{}{}}
\makeatother
\crefname{originallemma}{original lemma}{original lemmas}
\Crefname{originallemma}{Original lemma}{Original lemmas}
\makeatletter
\newlabel{lemma-trace-map-and-base-change}{{0B6J}{source}{}{sourcepointer-4}{}}
\newlabel{lemma-trace-map-and-base-change@cref}{{[originallemma][0][]0B6J}{[1][0][]source}{}{}{}}
\makeatother
\crefname{originallemma}{original lemma}{original lemmas}
\Crefname{originallemma}{Original lemma}{Original lemmas}
\makeatletter
\newlabel{lemma-flat-precompose-pus}{{0A9K}{source}{}{sourcepointer-5}{}}
\newlabel{lemma-flat-precompose-pus@cref}{{[originallemma][0][]0A9K}{[1][0][]source}{}{}{}}
\makeatother
\crefname{originallemma}{original lemma}{original lemmas}
\Crefname{originallemma}{Original lemma}{Original lemmas}
\makeatletter
\newlabel{categories-lemma-properties-2-cat-cats}{{003F}{source}{}{sourcepointer-6}{}}
\newlabel{categories-lemma-properties-2-cat-cats@cref}{{[originallemma][0][]003F}{[1][0][]source}{}{}{}}
\makeatother
\crefname{originaldefinition}{original definition}{original definitions}
\Crefname{originaldefinition}{Original definition}{Original definitions}
\makeatletter
\newlabel{categories-definition-horizontal-composition}{{003E}{source}{}{sourcepointer-7}{}}
\newlabel{categories-definition-horizontal-composition@cref}{{[originaldefinition][0][]003E}{[1][0][]source}{}{}{}}
\makeatother
\crefname{originallemma}{original lemma}{original lemmas}
\Crefname{originallemma}{Original lemma}{Original lemmas}
\makeatletter
\newlabel{categories-lemma-transformation-between-functors-and-adjoints}{{0B65}{source}{}{sourcepointer-8}{}}
\newlabel{categories-lemma-transformation-between-functors-and-adjoints@cref}{{[originallemma][0][]0B65}{[1][0][]source}{}{}{}}
\makeatother
\title{The unit of the adjunction commutes with base change through the comparison transformations}
\author{The Stacks Project Authors}
\date{stacks-project:duality.tex:local-snapshot-20261012; GNU FDL 1.2 or later}
\begin{document}
\maketitle
\noindent\textit{Editorial note.} All delivered mathematical text is the author's own, copied from the pinned Stacks Project chapter file duality.tex; no statement or proof step was rewritten, reordered or removed.\par\medskip
\noindent\textit{Editorial note.} Preamble provenance: the pinned chapter file opens with an input of the shared preamble and carries no preamble text of its own, so the project's single Stacks-wide preamble staged at sources/stacks-duality/preamble.tex is excerpted verbatim under the recipe fixed by LANE\_BRIEF 22.1, as three spans, lines 8-13 (verbatim package and the reference/slogan/history environments), lines 14-18 (Xy-pic setup) and lines 154-217 (the native theorem declarations and mathematical macros). Nothing was added, removed or redefined; the wrapper class and the amsmath/amssymb/amsthm/graphicx base packages are supplied by the assembler.\par\medskip
\noindent\textit{Editorial note.} The counted statement and proof contain no cross reference, no citation and no included file, so the delivered body is self contained.\par\medskip
\noindent\textit{Editorial note.} Canonical location of the counted proposition: Stacks tag 0B6K (sources/stacks-duality/pinned-original-source.tex line 1956).\par\medskip
\noindent\textit{Allowed standard prerequisites (editorial).}\begin{itemize}
\item The base change diagram and the comparison transformations alpha and beta. Same chapter, equation labelled equation-base-change-code, diagram labelled equation-base-change, equation labelled equation-base-change-map and Lemma labelled lemma-flat-precompose-pus, cited at proof lines 1978-1992. Fix the two transformations whose agreement is asserted.
\item The ordinary base change remark for pushforward. Cohomology, Remark labelled cohomology-remark-base-change, invoked in the statement. Supplies beta and its isomorphism property.
\item The dual trace map lemma. Same chapter, Lemma labelled lemma-trace-map-and-base-change, cited at proof lines 1978-1980. Provides the dual argument being imitated.
\item Coherence for compositions of functors and the transformation between adjoints. Categories, Lemma labelled categories-lemma-properties-2-cat-cats and Definition labelled categories-definition-horizontal-composition, cited at proof lines 2014-2018, and Categories, Lemma labelled categories-lemma-transformation-between-functors-and-adjoints, cited at proof lines 2020-2024. Give the commutativity of the squares and of the pentagon in the proof diagram.
\end{itemize}\medskip
\ifdefined\hypertarget\hypertarget{sourcepointer-1}{}\fi
\noindent\textit{Source reference.} Original equation 0A9J: sources/stacks-duality/pinned-original-source.tex, source lines 927--927. Same-chapter pointer into the author environment equation at source line 927; invoked at proof line None of the counted proof. It is retained only as a navigation pointer; every load-bearing step is inside the delivered span.\par\medskip
\ifdefined\hypertarget\hypertarget{sourcepointer-2}{}\fi
\noindent\textit{Source reference.} Original reference 08HY: sources/stacks-chow/tags.txt, source lines 10796--10796. Cross-chapter standard prerequisite located in the project tag list staged in the local snapshot; invoked at proof line None of the counted proof. It is retained only as a navigation pointer; every load-bearing step is inside the delivered span.\par\medskip
\ifdefined\hypertarget\hypertarget{sourcepointer-3}{}\fi
\noindent\textit{Source reference.} Original equation 0AA6: sources/stacks-duality/pinned-original-source.tex, source lines 1192--1192. Same-chapter pointer into the author environment equation at source line 1192; invoked at proof line None of the counted proof. It is retained only as a navigation pointer; every load-bearing step is inside the delivered span.\par\medskip
\ifdefined\hypertarget\hypertarget{sourcepointer-4}{}\fi
\noindent\textit{Source reference.} Original lemma 0B6J: sources/stacks-duality/pinned-original-source.tex, source lines 1867--1867. Same-chapter pointer into the author environment lemma at source line 1867; invoked at proof line 1979 of the counted proof. It is retained only as a navigation pointer; every load-bearing step is inside the delivered span.\par\medskip
\ifdefined\hypertarget\hypertarget{sourcepointer-5}{}\fi
\noindent\textit{Source reference.} Original lemma 0A9K: sources/stacks-duality/pinned-original-source.tex, source lines 957--957. Same-chapter pointer into the author environment lemma at source line 957; invoked at proof line 1986 of the counted proof. It is retained only as a navigation pointer; every load-bearing step is inside the delivered span.\par\medskip
\ifdefined\hypertarget\hypertarget{sourcepointer-6}{}\fi
\noindent\textit{Source reference.} Original lemma 003F: sources/stacks-chow/tags.txt, source lines 132--132. Cross-chapter standard prerequisite located in the project tag list staged in the local snapshot; invoked at proof line 2015 of the counted proof. It is retained only as a navigation pointer; every load-bearing step is inside the delivered span.\par\medskip
\ifdefined\hypertarget\hypertarget{sourcepointer-7}{}\fi
\noindent\textit{Source reference.} Original definition 003E: sources/stacks-chow/tags.txt, source lines 131--131. Cross-chapter standard prerequisite located in the project tag list staged in the local snapshot; invoked at proof line 2017 of the counted proof. It is retained only as a navigation pointer; every load-bearing step is inside the delivered span.\par\medskip
\ifdefined\hypertarget\hypertarget{sourcepointer-8}{}\fi
\noindent\textit{Source reference.} Original lemma 0B65: sources/stacks-chow/tags.txt, source lines 14058--14058. Cross-chapter standard prerequisite located in the project tag list staged in the local snapshot; invoked at proof line 2020 of the counted proof. It is retained only as a navigation pointer; every load-bearing step is inside the delivered span.\par\medskip
\begin{lemma}
\label{lemma-unit-and-base-change}
Suppose we have a diagram (\ref{equation-base-change}) where $f$ and $g$
are tor independent. Then the maps
$1 \star \eta_f : L(g')^* \to L(g')^* \circ a \circ Rf_*$ and
$\eta_{f'} \star 1 : L(g')^* \to a' \circ Rf'_* \circ L(g')^*$
agree via the base change maps
$\beta : Lg^* \circ Rf_* \to Rf'_* \circ L(g')^*$
(Cohomology, Remark \ref{cohomology-remark-base-change})
and $\alpha : L(g')^* \circ a \to a' \circ Lg^*$
(\ref{equation-base-change-map}).
More precisely, the diagram
$$
\xymatrix{
L(g')^* \ar[r]_-{1 \star \eta_f} \ar[d]_{\eta_{f'} \star 1} &
L(g')^* \circ a \circ Rf_* \ar[d]^\alpha \\
a' \circ Rf'_* \circ L(g')^* &
a' \circ Lg^* \circ Rf_* \ar[l]_-\beta
}
$$
of transformations of functors commutes.
\end{lemma}

\begin{proof}
This proof is dual to the proof of Lemma \ref{lemma-trace-map-and-base-change}.
In this proof we write $f_*$ for $Rf_*$ and $g^*$ for $Lg^*$ and we
drop $\star$ products with identities as one can figure out which ones
to add as long as the source and target of the transformation is known.
Recall that $\beta : g^* \circ f_* \to f'_* \circ (g')^*$ is an isomorphism
and that $\alpha$ is defined using
the isomorphism $\beta^\vee : g'_* \circ a' \to a \circ g_*$
which is the adjoint of $\beta$, see Lemma \ref{lemma-flat-precompose-pus}
and its proof. First we note that the left vertical arrow
of the diagram in the lemma is equal to the composition
$$
(g')^* \to (g')^* \circ g'_* \circ (g')^* \to
(g')^* \circ g'_* \circ a' \circ f'_* \circ (g')^* \to
a' \circ f'_* \circ (g')^*
$$
where the first arrow is the unit for $((g')^*, g'_*)$, the second arrow
is $\eta_{f'}$, and the third arrow is the counit for $((g')^*, g'_*)$.
This is a simple consequence of the fact that the composition
$(g')^* \to (g')^* \circ (g')_* \circ (g')^* \to (g')^*$
of unit and counit is the identity. Consider the diagram
$$
\xymatrix{
& (g')^* \circ a \circ f_* \ar[r] &
(g')^* \circ a \circ g_* \circ g^* \circ f_*
\ar[ld]_\beta \\
(g')^* \ar[ru]^{\eta_f} \ar[dd]_{\eta_{f'}} \ar[rd] &
(g')^* \circ a \circ g_* \circ f'_* \circ (g')^* &
(g')^* \circ g'_* \circ a' \circ g^* \circ f_*
\ar[u]_{\beta^\vee} \ar[ld]_\beta \ar[d] \\
& (g')^* \circ g'_* \circ a' \circ f'_* \circ (g')^*
\ar[ld] \ar[u]_{\beta^\vee} &
a' \circ g^* \circ f_* \ar[lld]^\beta \\
a' \circ f'_* \circ (g')^*
}
$$
In this diagram the two squares commute 
Categories, Lemma \ref{categories-lemma-properties-2-cat-cats}
or more simply the discussion preceding
Categories, Definition \ref{categories-definition-horizontal-composition}.
The triangle commutes by the discussion above. By the dual of
Categories, Lemma
\ref{categories-lemma-transformation-between-functors-and-adjoints}
the square
$$
\xymatrix{
\text{id} \ar[r] \ar[d] &
g'_* \circ a' \circ g^* \circ f_* \ar[d]^\beta \\
g'_* \circ a' \circ g^* \circ f_* \ar[r]^{\beta^\vee} &
a \circ g_* \circ f'_* \circ (g')^*
}
$$
commutes which implies the pentagon in the big diagram commutes.
Since $\beta$ and $\beta^\vee$ are isomorphisms, and since going on
the outside of the big diagram equals
$\beta \circ \alpha \circ \eta_f$ by definition this proves the lemma.
\end{proof}

\end{document}
